% Figura 3.3: Descomposición del precio de la acción en valor sin crecimiento y VAOC (Gordon-Shapiro y Tema 2 UGR)
\begin{tikzpicture}[
    font=\small\sffamily,
    lbl/.style={font=\scriptsize\sffamily, text=#1}
]
% Tres escenarios comparativos en barras apiladas: ROE < r (destruye), ROE = r (neutro), ROE > r (VAOC > 0)
% Eje base en y = 0
\draw[line width=0.9pt, color=dfinslate] (0.5, 0) -- (11.5, 0);
\draw[->, >=Latex, line width=0.9pt, color=dfinslate] (0.5, -1.3) -- (0.5, 5.4)
    node[above right, font=\footnotesize\sffamily\bfseries, color=dfinnavy] {Valor de la acción ($P_0$)};

% Referencia BPA1/r constante = 3.2 cm (ej. 50 u.m.)
\draw[dashed, color=dfinmuted, line width=0.8pt] (0.5, 3.2) -- (11.2, 3.2)
    node[right, font=\scriptsize\sffamily, color=dfinmuted] {$\dfrac{\text{BPA}_1}{r}$};

% Caso 1: Empresa vaca lechera / madura (b=0, ROE=r -> VAOC = 0)
\fill[dfinnavy!85] (1.3, 0) rectangle (3.3, 3.2);
\draw[color=dfinnavy, line width=0.9pt] (1.3, 0) rectangle (3.3, 3.2);
\node[font=\scriptsize\sffamily\bfseries, color=white, align=center] at (2.3, 1.6)
    {Valor sin\\[-1pt]crecimiento\\$\text{BPA}_1/r$};
\node[below=5pt, align=center, font=\scriptsize\sffamily, color=dfinslate] at (2.3, -1.3)
    {\textbf{Empresa madura}\\$b=0$ o $\text{ROE}=r$\\[-1pt]$\text{VAOC}=0$};
\node[above=3pt, font=\scriptsize\sffamily\bfseries, color=dfinnavy] at (2.3, 3.2) {$P_0 = \text{BPA}_1/r$};

% Caso 2: Empresa en crecimiento rentable (ROE > r, b > 0 -> VAOC > 0)
\fill[dfinnavy!85] (4.7, 0) rectangle (6.7, 3.2);
\fill[dfinemerald!85] (4.7, 3.2) rectangle (6.7, 4.9);
\draw[color=dfinnavy, line width=0.9pt] (4.7, 0) rectangle (6.7, 3.2);
\draw[color=dfinemerald, line width=0.9pt] (4.7, 3.2) rectangle (6.7, 4.9);
\node[font=\scriptsize\sffamily\bfseries, color=white, align=center] at (5.7, 1.6)
    {$\text{BPA}_1/r$};
\node[font=\scriptsize\sffamily\bfseries, color=white, align=center] at (5.7, 4.05)
    {$+\text{VAOC} > 0$\\[-1pt]{\tiny ($\text{ROE} > r$)}};
\node[below=5pt, align=center, font=\scriptsize\sffamily, color=dfinslate] at (5.7, -1.3)
    {\textbf{Valor de crecimiento}\\$\text{ROE} > r$, retiene $b>0$\\[-1pt]Crea riqueza};
\node[above=3pt, font=\scriptsize\sffamily\bfseries, color=dfinemerald] at (5.7, 4.9) {$P_0 > \text{BPA}_1/r$};

% Caso 3: Crecimiento destructor de valor (ROE < r, b > 0 -> VAOC < 0)
\fill[dfinnavy!85] (8.1, 0) rectangle (10.1, 3.2);
\fill[dfinburg!82] (8.1, -1.0) rectangle (10.1, 0);
\draw[color=dfinnavy, line width=0.9pt] (8.1, 0) rectangle (10.1, 3.2);
\draw[color=dfinburg, line width=0.9pt] (8.1, -1.0) rectangle (10.1, 0);
\node[font=\scriptsize\sffamily\bfseries, color=white, align=center] at (9.1, 1.6)
    {$\text{BPA}_1/r$};
\node[font=\scriptsize\sffamily\bfseries, color=white, align=center] at (9.1, -0.5)
    {$\text{VAOC} < 0$};
\node[below=5pt, align=center, font=\scriptsize\sffamily, color=dfinslate] at (9.1, -1.3)
    {\textbf{Sobreinversión}\\$\text{ROE} < r$, retiene $b>0$\\[-1pt]Destruye valor};
% Llave indicando P0 neto inferior
\draw[decorate, decoration={brace, amplitude=5pt, mirror}, color=dfinburg, line width=1pt]
    (10.25, 1.0) -- (10.25, 3.2) node[midway, right=6pt, font=\scriptsize\sffamily, color=dfinburg, align=left]
    {$P_0$ neto\\inferior a\\$\text{BPA}_1/r$};

\end{tikzpicture}
